\documentclass[11pt]{article}

\usepackage[utf8]{inputenc}
\usepackage[T1]{fontenc}
\usepackage{lmodern}
\usepackage{geometry}
\usepackage{amsmath,amssymb,amsfonts,amsthm,mathtools}
\usepackage{bm}
\usepackage{enumitem}
\usepackage{microtype}
\usepackage{hyperref}
\usepackage{titlesec}
\usepackage{booktabs}
\usepackage{array}
\usepackage{setspace}
\usepackage{mathrsfs}
\usepackage{physics}

\geometry{margin=1in}
\hypersetup{colorlinks=true, linkcolor=black, urlcolor=blue, citecolor=black}
\setlist[enumerate]{topsep=0.4em,itemsep=0.35em}
\setlist[itemize]{topsep=0.3em,itemsep=0.25em}
\titleformat{\section}{\large\bfseries}{\thesection.}{0.5em}{}
\titleformat{\subsection}{\normalsize\bfseries}{\thesubsection.}{0.5em}{}
\setstretch{1.06}

\newcommand{\Aether}{\AE{}ther}
\newcommand{\Aflow}{\AE{}ther-flow}
\newcommand{\TheoryName}{\Aflow{} Interpretation of Relativity}
\newcommand{\TheoryProgram}{\Aether{} / \Aflow{} framework}
\newcommand{\CompletedMetric}{g^{\sharp}}
\newcommand{\SubPreSectionCore}{\mathfrak U^{\mathrm{sec}}}
\newcommand{\SubSectionPackage}{\mathfrak V^{\mathrm{subsec}}}

\newtheorem{definition}{Definition}
\newtheorem{proposition}{Proposition}
\newtheorem{remark}{Remark}
\newtheorem{corollary}{Corollary}
\newtheorem{theorem}{Theorem}

\title{\textbf{The \TheoryName{}}\\[0.4em]
\large Primitive-Reservoir Observer-Reduction\\
\large Exact-Shell-Transport-Section-Generating Substrate\\
\large Sub-Section Dynamics Next Benchmark Criterion}
\author{}
\date{}

\begin{document}

\maketitle

\begin{abstract}
The active primitive-reservoir observer-reduction line has now moved below the observer-relevant sub-pre-core exact-shell transport-section core itself. The transport-image collapse theorem already spent the sharper internal compression move available inside the sub-pre-core image presentation, and the later sub-section-dynamics theorem then removed the section core \(\SubPreSectionCore\) as the primitive stopping point by deriving it canonically from a deeper exact-shell-transport-section-generating substrate sub-section dynamics package \(\SubSectionPackage\).

This note records the routing consequence of that theorem. Another same-output deeper-origin relay that merely preserves the generated section core \(\SubPreSectionCore\) no longer counts as the honest default next benchmark-facing manuscript. Neither does a reopening of the larger transport-image, full sub-pre-core, or already-screened pre-core packages as if they were again independently live. The next honest benchmark-facing move must therefore either derive or justify the deeper sub-section dynamics package \(\SubSectionPackage\) from still more primitive substrate structure in a way that changes benchmark-facing status, or prove a genuinely smaller theorem-level collapse strictly inside that deeper package.

The benchmark boundary remains fixed throughout. The one operative metric is still exactly \(\CompletedMetric\), the ordinary matter route is unchanged, there is no second operative metric, the low-energy matter law is not enlarged, and no stronger promotion language is licensed. The positive result of this note is therefore routing discipline at the current sub-section-dynamics frontier.
\end{abstract}

\section{Introduction}

The active derivational program of \TheoryName{} is no longer at a stage where the observer-relevant sub-pre-core exact-shell transport-section core should be treated as the default unexplained stopping point. The earlier image-core collapse theorem had already shown that the larger transport-image presentation carried no additional benchmark-relevant freedom beyond the smaller section core \(\SubPreSectionCore\). The later sub-section-dynamics theorem then spent the next honest move available there by showing that even this section core is canonically generated by a deeper exact-shell-transport-section-generating substrate sub-section dynamics package \(\SubSectionPackage\).

The present note records the routing consequence of that deeper theorem. It does not reopen the sub-section mathematics, and it does not claim a completed first-principles substrate derivation of gravity. It records only which kind of next manuscript would now count as an honest benchmark-facing gain below the current sub-section-dynamics frontier.

\paragraph{Framework and claim boundary.}
\TheoryProgram{} denotes the broader \Aether{} / \Aflow{} framework built on the claims that the \Aether{} is the underlying four-dimensional substrate of reality, the \Aflow{} is its intrinsic ordered motion, observed three-dimensional space is the local experiential slice of that deeper substrate, S-time is the experienced order of change arising from matter, light, and the \Aflow{}, and observed expansion is the three-dimensional appearance of deeper four-dimensional ordered motion. Within that framework, \TheoryName{} denotes the exact-closure theory in which the effective gravitational dynamics are adopted to be exactly Einsteinian with universal matter coupling. In the active sequence, \emph{adoption} means use of that established relativistic dynamics without claiming substrate derivation, while \emph{derivation} is reserved for a first-principles recovery from explicit substrate variables.

The present note stays strictly inside that boundary. It records only which kind of next manuscript would now count as an honest benchmark-facing gain at the current sub-section-dynamics frontier.

\section{Fixed Inputs}

\begin{definition}[Sub-section-dynamics frontier inputs]
The criterion recorded here uses the following fixed inputs.
\begin{enumerate}
    \item The benchmark exact-closure package, which fixes the public one-metric GR target of the repository.
    \item The benchmark gatekeeping layer, including the recorded safeguard that blocks same-output deeper-origin relays from counting as new front-line progress once the live observer-relevant datum is unchanged.
    \item The transport-image collapse theorem, which compresses the sub-pre-core image presentation to the smaller observer-relevant sub-pre-core exact-shell transport-section core \(\SubPreSectionCore\).
    \item The later sub-section-dynamics theorem, which proves that the section core \(\SubPreSectionCore\) itself is canonically generated by the deeper exact-shell-transport-section-generating substrate sub-section dynamics package \(\SubSectionPackage\).
\end{enumerate}
\end{definition}

\begin{remark}
The present note does not replace those manuscripts. It records the routing consequence of reading them together under the active safeguard against recursive depth drift.
\end{remark}

\section{Why the Sub-Section-Dynamics Theorem Is the Frontier}

\begin{definition}[Benchmark-neutral deeper relay at sub-section scope]
A \emph{benchmark-neutral deeper relay at sub-section scope} is a manuscript that preserves the same benchmark one-metric output, the same ordinary matter route, the same generated section core \(\SubPreSectionCore\), and the same deeper sub-section package \(\SubSectionPackage\) while merely relocating the unexplained substrate layer deeper, restating the larger transport-image or full sub-pre-core presentations, reopening already-screened pre-core package data, or repackaging the same sub-section package without a new compression, necessity, uniqueness, or comparably sharp routing gain.
\end{definition}

\begin{proposition}[The sub-section-dynamics theorem is the live frontier]
At the current stage of the primitive-reservoir observer-reduction line, the theorem deriving the observer-relevant sub-pre-core exact-shell transport-section core from exact-shell-transport-section-generating substrate sub-section dynamics remains the live benchmark-facing frontier.
\end{proposition}

\begin{proof}
That theorem changes benchmark-facing status at current scope because it removes the section core \(\SubPreSectionCore\) itself as the unexplained stopping point. The live burden now lies at the deeper sub-section dynamics package \(\SubSectionPackage\), not at the previously exposed section core. By contrast, a later manuscript that preserves the same benchmark package, the same generated section core, and the same deeper package while merely moving the unexplained layer deeper does not alter benchmark-facing status unless it adds a comparably sharp gain. Therefore the live frontier is now the sub-section-dynamics theorem rather than the earlier transport-image collapse theorem or any later same-output relay that leaves the deeper package unchanged.
\end{proof}

\begin{corollary}[Status of the prior transport-image collapse theorem]
The theorem collapsing the observer-relevant sub-pre-core exact-shell transport-image core to the smaller transport-section core remains preserved benchmark-facing main-line history, but under the current routing safeguard it is no longer by itself the default current frontier.
\end{corollary}

\begin{proof}
That theorem matters because it proved that the larger sub-pre-core image presentation carries no additional benchmark-relevant freedom beyond the smaller section core \(\SubPreSectionCore\). But the later sub-section-dynamics theorem shows that even \(\SubPreSectionCore\) is no longer the present unexplained stopping point. The earlier theorem is preserved history and handoff, while the later theorem defines the current frontier.
\end{proof}

\section{The Next Benchmark Criterion}

\begin{definition}[Admissible next benchmark-facing manuscript at sub-section scope]
At the current sub-section-dynamics frontier, a manuscript counts as an admissible next benchmark-facing move only if it does at least one of the following:
\begin{enumerate}
    \item derives or justifies the exact-shell-transport-section-generating substrate sub-section dynamics package \(\SubSectionPackage\) from still more primitive substrate structure in a way that does more than merely relocate the unexplained layer deeper;
    \item proves a genuinely smaller theorem-level collapse strictly inside \(\SubSectionPackage\);
    \item does either of the previous two while preserving the same benchmark one-metric package, the same ordinary matter route, and the same claim boundary.
\end{enumerate}
\end{definition}

\begin{theorem}[Next honest benchmark-facing task at sub-section scope]
The next honest benchmark-facing manuscript below the current sub-section-dynamics frontier must either derive or justify the exact-shell-transport-section-generating substrate sub-section dynamics from still more primitive substrate structure in a benchmark-relevant way, or else prove a genuinely smaller collapse strictly inside that deeper package. A manuscript that only repeats the same generated section core through another deeper relay, re-expands the discussion back to the larger transport-image or sub-pre-core packages without such a new gain, or invokes a quantum branch without doing the same benchmark-facing work does not satisfy the criterion.
\end{theorem}

\begin{proof}
By Proposition 1, the live frontier already lies at the deeper sub-section package \(\SubSectionPackage\) rather than at the generated section core \(\SubPreSectionCore\) or the larger transport-image presentation. Therefore a subsequent manuscript must improve on that status rather than merely accompany it. A deeper-origin manuscript can do so only if it changes benchmark-facing status by more than depth alone. If it simply reproduces the same generated section core through another relay, or merely re-expands the already-screened larger packages, the routing rule classifies it as benchmark neutral. If it invokes an explicitly quantum package without deriving the same deeper sub-section dynamics, collapsing them further, or closing a benchmark derivation gate, it likewise leaves the live burden unchanged. The remaining honest alternatives are therefore exactly those stated in the definition: either a more primitive derivation or justification of the sub-section package itself, or a sharper internal collapse of that package.
\end{proof}

\section{What the Next Move Should Not Be}

\begin{proposition}[Screened next moves]
At the current sub-section-dynamics frontier, none of the following is the default next benchmark-facing move:
\begin{enumerate}
    \item another same-output deeper-origin relay that merely preserves the generated section core \(\SubPreSectionCore\);
    \item another re-expansion back to the larger transport-image, full sub-pre-core, or already-screened pre-core presentations as if they were again independently live burdens;
    \item a quantum branch that does not derive or justify the same deeper sub-section dynamics, collapse them further, or close a benchmark derivation gate in a genuinely new way;
    \item a manuscript that introduces a second operative metric, enlarges the ordinary matter law, or uses stronger promotion language.
\end{enumerate}
\end{proposition}

\begin{proof}
Items (1) and (2) fail the preceding criterion because they do not directly address the current sub-section-level burden. They revisit a same-output relay or a previously screened larger presentation rather than removing the remaining freedom in \(\SubSectionPackage\) or proving a smaller internal datum there. Item (3) likewise does not directly address the live burden at current scope. It changes vocabulary without changing the benchmark-facing task unless it derives the same deeper package, sharpens it internally, or settles a derivation gate. Item (4) violates the fixed claim boundary of the benchmark package and therefore cannot count as a disciplined next step on the active line. The benchmark package remains the exact one-metric GR package already fixed by adoption, so any admissible next manuscript must preserve that boundary.
\end{proof}

\begin{remark}
This screening statement is not a denial that those deeper or alternative manuscripts may matter strategically. It is a routing statement: they are not the default way to spend the next benchmark-facing move from the current sub-section-dynamics frontier unless they do the same benchmark-facing work named above.
\end{remark}

\section{Conclusion}

The positive result of this note is routing discipline at the current sub-section-dynamics frontier. The sub-section-dynamics theorem remains the live benchmark-facing gain because it already removes the observer-relevant sub-pre-core exact-shell transport-section core as the unexplained stopping point on the active line. The earlier transport-image collapse theorem remains preserved main-line history, but under the recursive-depth safeguard it is no longer itself the default current frontier.

The scope remains narrow and honest. The benchmark package is unchanged: the one operative metric is still exactly \(\CompletedMetric\), the ordinary matter route is unchanged, there is no second operative metric, and the low-energy matter law is not enlarged. Nothing here upgrades adoption to derivation or licenses stronger promotion language.

The next open burden is therefore clear. The next benchmark-facing manuscript should derive or justify the exact-shell-transport-section-generating substrate sub-section dynamics from still more primitive substrate structure in a way that changes benchmark-facing status, unless the sharper honest gain available instead is a genuinely smaller theorem-level collapse strictly inside that deeper package. Another same-output deeper relay that merely preserves the generated section core, a reopening of the larger screened packages, or a quantum branch that does not do that same benchmark-facing work is not the clean next manuscript target at current scope.

\end{document}
